\part{New in version 0.3}
\chapter{New modules in version 0.3}
\label{ch:new-modules}

Version 0.3 adds five layers to the public surface: a relaxed-path
NEVPT2 nuclear gradient, a coupled CASPT2 contraction, a valence-level
spin-orbit mean field for effective-core potentials, a RIJONX
self-consistent-field driver, and a scalar density-fitted X2C CASSCF
layer.  Each is additive: no existing signature changed, and the
release core described in Chapter~\ref{ch:differentiable-chain} is
untouched.  This chapter states, for each layer, what the entry point
computes, what the caller must supply, and where the implementation
stops.  The RI factor conventions reused by two of the layers are those
of Chapter~\ref{ch:ri-df}.

\begin{center}
\small
\begin{tabularx}{\textwidth}{>{\raggedright\arraybackslash}p{0.33\textwidth}>{\raggedright\arraybackslash}p{0.36\textwidth}X}
\toprule
Module & Main entry point & Layer \\
\midrule
\texttt{driver.nevpt2\_zvector} &
  \texttt{nevpt2\_relaxed\_gradient\_zvector} &
  response on top of the differentiated chain \\
\texttt{core.caspt2} &
  \texttt{caspt2\_contracted\_e2\_coupled} &
  value (NumPy) \\
\texttt{core.somf\_ecp} &
  \texttt{valence\_somf\_matrix} &
  value (NumPy) \\
\texttt{driver.ri\_scf}, \texttt{driver.ri\_scf\_driver} &
  \texttt{rijonx\_scf}, \texttt{rijonx\_scf\_gradient} &
  SCF (pyscfad) \\
\texttt{core.x2c\_casscf\_df} &
  \texttt{df\_x2c\_casscf} &
  value and analytic gradient \\
\bottomrule
\end{tabularx}
\end{center}

\section{Relaxed NEVPT2 nuclear gradients}
\label{sec:new-nevpt2-zvector}

\texttt{driver.nevpt2.nevpt2\_nuclear\_gradient} differentiates the
frozen-theta energy: the SCF reference and the CI vector are rebuilt at
every displaced geometry, but the orbital-rotation parameters
$\theta$ stay at their stationary values.  Along the physical path
$\theta$ follows the CASSCF stationarity condition, and NEVPT2 is not
variational in the orbitals, so the frozen partial is not the relaxed
derivative.  The correction is the Z-vector response

\begin{equation}
  g_{\mathrm{relaxed}} = g_{\mathrm{frozen}} - M^{T} H_{\theta\theta}^{-1} s,
  \qquad
  H_{\theta\theta} = \frac{\partial^2 E_{\mathrm{CASSCF}}}
    {\partial\theta^{2}}, \quad
  M = \frac{\partial^2 E_{\mathrm{CASSCF}}}
    {\partial\theta\,\partial R}, \quad
  s = \frac{\partial e_2}{\partial\theta},
  \label{eq:new-zvector}
\end{equation}

with $e_2 = E_{\mathrm{NEVPT2}} - E_{\mathrm{CASSCF}}$.  The vector
$z = H_{\theta\theta}^{-1} s$ is the orbital Z-vector: one linear solve
per state serves every Cartesian direction.  The CI coordinates need no
separate multiplier in this reduced space, because both energy
functions solve the CI eigenvalue problem internally; the CI block of
the full Hessian is replaced by its Schur complement, which is
$H_{\theta\theta}$ as written.

\begin{center}
\small
\begin{tabularx}{\textwidth}{>{\raggedright\arraybackslash}p{0.43\textwidth}>{\raggedright\arraybackslash}X}
\toprule
Entry point & Returns and purpose \\
\midrule
\texttt{make\_nevpt2\_joint\_fn\_si\_exact(...)} &
  $E(\mathrm{mol},\theta)$ with the exact near-zero \texttt{Si}
  channels; same contract as
  \texttt{driver.nevpt2.make\_nevpt2\_joint\_fn}.  Takes the space
  definition, the orbital pairs, the SCF builder, \texttt{e2\_scale},
  and optionally \texttt{align\_ref} and \texttt{align\_bounds}. \\
\texttt{align\_reference\_blocks(...)} &
  Differentiable blockwise alignment of the SCF coefficients onto a
  fixed reference; arguments \texttt{(C, S, C\_ref, bounds)}, where
  \texttt{bounds} is a list of \texttt{(lo, hi)} orbital ranges. \\
\texttt{orbital\_zvector(...)} &
  Solves $H_{\theta\theta} z = s$ by truncated pseudo-inverse; arguments
  \texttt{(H\_tt, source, rcond=None, symmetrize=True)}; returns
  \texttt{(z, info)} with the spectrum, condition number and residual. \\
\texttt{lagrangian\_relaxed\_gradient(...)} &
  The contraction $g_{\mathrm{frozen}} - M^{T} z$; arguments
  \texttt{(g\_frozen, z, M)}. \\
\texttt{nevpt2\_relaxed\_gradient\_zvector(...)} &
  The assembled relaxed gradient; returns \texttt{(gradient, info)}.
  Required arguments: \texttt{mol}, \texttt{theta}, \texttt{E\_casscf},
  \texttt{E\_nevpt2}; optional \texttt{rcond}, \texttt{H\_tt},
  \texttt{M}, \texttt{source}, \texttt{g\_frozen}. \\
\texttt{si\_energy\_from\_invariants(...)} &
  The \texttt{Si} channel sum over \texttt{(norm, ener, diff)};
  \texttt{mode='masked'} reproduces the release value mask for
  diagnostics only. \\
\bottomrule
\end{tabularx}
\end{center}

Three settings carry the accuracy of the response.

\begin{itemize}
\item \emph{Reference convergence.}  The reference SCF is converged with
  \texttt{conv\_tol} of $10^{-13}$ and \texttt{conv\_tol\_grad} of
  $10^{-10}$ at every geometry.  The response operator is a second
  derivative of the same chain, so its accuracy is bounded by the
  orbital-gradient convergence of the reference, not by the energy
  convergence alone.
\item \emph{Blockwise alignment.}  The SCF solution is defined only up
  to orthogonal mixing inside the core, active and virtual blocks, and
  two neighbouring geometries can return different representatives (a
  phase flip, or an ordering change inside a near-degenerate pair).
  Because $\theta$ is defined relative to that representative, an
  unaligned reference makes the same $\theta$ mean different orbitals:
  on H$_2$O/6-31G the measured fixed-theta energy slope was
  $-5.8$~Eh/bohr against $-0.0528$~Eh/bohr in the aligned gauge.  Pass
  the base-geometry coefficients as \texttt{align\_ref} and the block
  bounds as \texttt{align\_bounds}; the map is the identity at the base
  geometry.
\item \emph{Exact near-zero \texttt{Si} channels.}  The release core
  evaluates the \texttt{Si} class as
  $-\sum_i n_i/(d_i + h_i/n_i)$ with a value mask $|n_i| > 10^{-12}$,
  where $n_i$ is the channel norm, $d_i$ its zeroth-order denominator
  and $h_i$ the core-coupling term.  Symmetry can drive an individual
  $n_i$ to $10^{-13}$: dropping such a channel changes the energy by
  about $10^{-13}$~Eh but its derivative by order $0.1$--$1$, because
  the mask is a discontinuity in the energy, not a removable
  singularity.  The module evaluates the same algebra as
  $-\sum_i n_i^2/(d_i n_i + h_i)$ with the guard moved to the
  denominator.  Only the \texttt{Si} class is treated this way; the
  other four contraction classes keep the release masks.
\end{itemize}

The validation covers closed-shell, frozen-core H$_2$O/6-31G and
LiH/6-31G, CAS(2,2), against finite differences of the relaxed path.
Open-shell, ECP, f-block and large-basis systems are outside the
validated range.  A minimal call is

\begin{lstlisting}[language=Python]
import jax.numpy as jnp
from pyscfad import scf

from diffmrc.core.orblib import orbital_pairs
from diffmrc.driver.nevpt2_zvector import (
    make_nevpt2_joint_fn_si_exact, nevpt2_relaxed_gradient_zvector)

ncore, nact, nele = 1, 2, 2
nmo = mol.nao_nr()
pairs = orbital_pairs(ncore, nact, nmo, include_cv=True)

def builder(m):
    mf = scf.RHF(m)
    mf.conv_tol = 1e-13
    mf.conv_tol_grad = 1e-10
    return mf

mf = builder(mol)
mf.kernel()
bounds = [(0, ncore), (ncore, ncore + nact), (ncore + nact, nmo)]
ref = jnp.asarray(mf.mo_coeff)

# theta: CASSCF stationary parameters at this geometry
E_casscf = make_nevpt2_joint_fn_si_exact(
    ncore, nact, nele, nmo, pairs, builder, e2_scale=0.0,
    align_ref=ref, align_bounds=bounds)
E_nevpt2 = make_nevpt2_joint_fn_si_exact(
    ncore, nact, nele, nmo, pairs, builder, e2_scale=1.0,
    align_ref=ref, align_bounds=bounds)

grad, info = nevpt2_relaxed_gradient_zvector(
    mol, theta, E_casscf, E_nevpt2, rcond=1e-7)
\end{lstlisting}

The cutoff \texttt{rcond} is a relative cutoff on the pseudo-inverse.
The H$_2$O CASSCF Hessian carries one soft mode (eigenvalue
$3.8\times10^{-6}$, condition number $2.3\times10^{7}$) whose response
is not part of the correlation source; $10^{-7}$ cuts it, while a looser
cutoff keeps it and amplifies roundoff.

\section{Coupled CASPT2}
\label{sec:new-caspt2-coupled}

\texttt{core.caspt2} has two contraction schemes of the same eight-class
internally contracted basis.  The diagonal path
(\texttt{caspt2\_contracted\_e2}) evaluates each contracted function
with its own denominator,

\begin{equation}
  e_2 = -\frac{W^2}{A - E_0 S},
  \label{eq:new-caspt2-diagonal}
\end{equation}

which is the CASPT2-D approximation.  The coupled path replaces the
diagonal zeroth-order Hamiltonian by the state-specific generalized
Fock operator $H_0 = \sum_{pq} f_{pq} E_{pq}$.  Its off-diagonal blocks
(inactive--active, active--virtual, inactive--virtual) connect
determinants of different excitation classes, and the amplitudes solve
the projected linear system

\begin{equation}
  (A - E_0 S)\,x = -W, \qquad
  A_{pq} = \langle K_p | H_0 | K_q \rangle, \quad
  S_{pq} = \langle K_p | K_q \rangle, \quad
  W_p = \langle K_p | H | 0 \rangle,
  \label{eq:new-caspt2-coupled}
\end{equation}

with $e_2 = \sum_p W_p x_p$ and the class decomposition
$e_{2,\mathrm{class}} = \sum_{p \in \mathrm{class}} W_p x_p$.  The
overlap matrix $S$ is regularly linearly dependent; it is diagonalized
and eigenvectors with eigenvalue below
$\texttt{rtol} \times \max(1, \lambda_{\max})$ are dropped.

\begin{center}
\small
\begin{tabularx}{\textwidth}{>{\raggedright\arraybackslash}p{0.40\textwidth}>{\raggedright\arraybackslash}X}
\toprule
Entry point & Returns and purpose \\
\midrule
\texttt{caspt2\_contracted\_e2\_coupled(...)} &
  \texttt{(e2\_total, breakdown)}, plus the assembled \texttt{S},
  \texttt{A}, \texttt{W}, \texttt{E0}, \texttt{x} and the class labels
  when \texttt{return\_system=True}.  Arguments: the MO integrals, the
  CI vector and determinant bitmaps, \texttt{ncore}, \texttt{nact},
  \texttt{nele}; optional \texttt{fock} and \texttt{rtol}.  With
  \texttt{fock} omitted, the state-specific generalized Fock is built
  from the CI vector and the integrals. \\
\texttt{caspt2\_contracted\_e2}, \texttt{caspt2\_contracted\_e2\_jax} &
  The diagonal-Fock path and its differentiable twin; the latter is the
  jax chain that can be traced in the CI coefficients. \\
\bottomrule
\end{tabularx}
\end{center}

The relative weight of the \texttt{VM} class follows the reference
implementation: with the two orderings written cross and direct, the
Molcas normalisation is $V_P = \text{cross} + \text{direct}$ and
$V_M = \sqrt{3}\,(\text{cross} - \text{direct})$.  The factor
$\sqrt{3}$ multiplies the whole \texttt{VM} function and cancels in
$W^2/(A - E_0 S)$, so the operator content carried here is
$\text{cross} - \text{direct}$, that is, a relative weight of $-1$
between the two orderings.  Both components enter the energy
independently for off-diagonal pairs and coincide for diagonal pairs.
Version 0.3 corrects the previous convention for this component.

The coupled total against the reference CASPT2 class table is $1.0081$
on H$_2$O/6-31G and $1.0100$ on CO/6-31G, both CAS(2,2).  The path is a
NumPy value layer: it has no jax twin, and the differentiable CASPT2
route remains the diagonal-Fock
contraction \texttt{caspt2\_contracted\_e2\_jax} at the CASPT2-D level.
The cost is $n_{\mathrm{fn}}^2$ function-pair contractions plus one
$H_0$ application per function, intended for small and medium active
spaces.

\begin{lstlisting}[language=Python]
from diffmrc.core.caspt2 import caspt2_contracted_e2_coupled

# h1e_mo, eri_mo: MO integrals in chemist notation, from the CASSCF
# orbitals; ci: reference CI vector; dets: determinant bitmaps.
e2, breakdown = caspt2_contracted_e2_coupled(
    h1e_mo, eri_mo, ci, dets, ncore=ncore, nact=nact, nele=nele)

# With an explicit state-specific generalized Fock matrix, and the
# assembled system for diagnostics:
e2, breakdown, system = caspt2_contracted_e2_coupled(
    h1e_mo, eri_mo, ci, dets, ncore=ncore, nact=nact, nele=nele,
    fock=fock, return_system=True)
\end{lstlisting}

\section{Valence-level SOMF for effective-core potentials}
\label{sec:new-somf-ecp}

The all-electron SOMF operator of \texttt{core.somf} uses the bare
nuclear charges and the full electron density.  With an effective core
potential the core electrons are absent, and the spin-orbit part of the
ECP replaces their contribution; adding the all-electron formula
unchanged would use the ECP effective charge in
\texttt{int1e\_pnucxp} and a valence-only density.  The module defines
the valence-level operator

\begin{equation}
  h^{\mathrm{SOMF\text{-}ECP}}_k
  = \frac{\alpha^{2}}{2}\left[h^{\mathrm{nuc}}_k + h^{2e}_k[P]\right]
    - V^{\mathrm{ECP}}_k,
  \label{eq:new-somf-ecp}
\end{equation}

with $h^{\mathrm{nuc}}_k$ the nuclear one-electron term on centers
without a spin-orbit ECP, $h^{2e}_k[P]$ the valence
$J - \tfrac{3}{2}K$ mean field of the spin-summed density $P$
($\mathrm{tr}(PS) = \texttt{mol.nelectron}$), and $V^{\mathrm{ECP}}_k$
the three-component \texttt{ECPso} integral.  On a center carrying a
spin-orbit ECP the ECP term replaces the nuclear term: the ECP is
fitted to reproduce the all-electron spin-orbit splitting of the
valence configuration and already contains the nuclear one-electron
term and the core-valence screening for that center.  The
Breit--Pauli terms carry the factor $\alpha^2/2$; \texttt{ECPso} is
already an effective operator in energy units and is subtracted.

\begin{center}
\small
\begin{tabularx}{\textwidth}{>{\raggedright\arraybackslash}p{0.40\textwidth}>{\raggedright\arraybackslash}X}
\toprule
Entry point & Returns and purpose \\
\midrule
\texttt{valence\_somf\_matrix(...)} &
  The operator, shape \texttt{(3, nao, nao)}.  Arguments:
  \texttt{(mol, P)}, then \texttt{nuclear='auto'}, \texttt{ecp\_term},
  \texttt{contract}, \texttt{memory\_limit}, \texttt{convention='i'},
  \texttt{check}.  \texttt{nuclear} selects which centers keep the
  nuclear term (\texttt{'auto'} excludes spin-orbit ECP centers);
  \texttt{convention='i'} returns
  $H_{\mathrm{SO}} = i\sum_k h_k s_k$. \\
\texttt{ecp\_soc(mol)} &
  The three-component \texttt{ECPso} integral, shape
  \texttt{(3, nao, nao)}, with no $\alpha^2/2$ factor. \\
\texttt{nuclear\_soc(mol, atoms=None)} &
  The nuclear one-electron SOC in the \texttt{int1e\_pnucxp}
  convention, optionally restricted to a list of centers. \\
\texttt{so\_ecp\_centers(mol)}, \texttt{nuclear\_centers(mol)} &
  Atom indices that carry a spin-orbit ECP, and those that do not. \\
\texttt{somf\_density\_response(...)} &
  The density-response matrix $M_{rs} = \sum_{pq} D_{pq}
  [g_{pqrs} - \tfrac{3}{2} g_{prsq} - \tfrac{3}{2} g_{sqpr}]$ for
  \texttt{(mol, D)}; only \texttt{contract='explicit'} is implemented. \\
\texttt{somf\_definition(mol)} &
  A dictionary recording which centers contribute which term and the
  approximations that are not fixed by the ECP. \\
\bottomrule
\end{tabularx}
\end{center}

The two-electron kernel is not symmetric under the pair exchange
$(pq) \leftrightarrow (rs)$.  Consequently the derivative of
$\mathrm{tr}(D\,h^{2e}[P])$ with respect to $P$ is not
$h^{2e}[D]$; a gradient of an SOC matrix element must use
\texttt{somf\_density\_response}, or the equivalent bilinear form.
The value layer is NumPy; orbital and CI response are the caller's
responsibility, and only the explicit contraction is implemented in the
response kernel.

There is no unique first-principles decomposition of an ECP into one-
and two-electron spin-orbit parts: the ECP is fitted to atomic
reference states, so part of the valence--valence mean field is already
folded into its parameters.  Adding the explicit valence
$J - \tfrac{3}{2}K$ term is therefore an approximation with a
double-counting error that cannot be removed without the ECP generation
data.  The operator is a prototype for relative comparisons and is not
a production absolute-SOC method.  The pieces are exposed independently
(\texttt{nuclear=}, \texttt{ecp\_term=}) so that variants can be
compared and the error bounded.

\begin{lstlisting}[language=Python]
from pyscf import gto

from diffmrc.core import somf_ecp

mol = gto.M(atom='Yb 0 0 0; F 0 0 3.6',
            basis={'Yb': 'stuttgart_rsc', 'F': '6-31g'},
            ecp={'Yb': 'ecpds28mwbso'}, unit='Bohr', verbose=0)

# P: spin-summed density of the electrons present in the calculation,
# trace(P S) = mol.nelectron.  D: the same density, or the transition
# 1-RDM for a matrix element between two scalar states.
h_soc = somf_ecp.valence_somf_matrix(mol, P, convention='i')
m_resp = somf_ecp.somf_density_response(mol, D)
\end{lstlisting}

On the YbF test molecule the \texttt{ECPso} norm on Yb is $17.04$ per
component, while the raw nuclear operator norm on the same center is of
order $10^{4}$; adding both would be dominated by an operator that the
ECP was built to replace.

\section{RIJONX SCF}
\label{sec:new-rijonx}

The reference SCF of the multireference chain still formed the
four-center Coulomb tensor through the dense integral path.  The RIJONX
route fits the Coulomb term and keeps the exchange exact:

\begin{equation}
  E = E_{\mathrm{nuc}} + \mathrm{Tr}(Dh)
    + \tfrac{1}{2}\mathrm{Tr}(D J_{\mathrm{RI}})
    - \tfrac{1}{4}\mathrm{Tr}(D K_{\mathrm{exact}}),
  \qquad
  J_{\mathrm{RI}} = W D W^{T},
  \label{eq:new-rijonx}
\end{equation}

with $W = B S^{-1/2}$ the metric-folded three-center tensor of
Chapter~\ref{ch:ri-df}.  The exchange branch is the dense four-center
integral, so the method is RIJONX: not RI-JK, which has no analytic
gradient here, and not RIJCOSX, whose fixed-grid prototype is outside
the production surface.

\begin{center}
\small
\begin{tabularx}{\textwidth}{>{\raggedright\arraybackslash}p{0.40\textwidth}>{\raggedright\arraybackslash}X}
\toprule
Entry point & Returns and purpose \\
\midrule
\texttt{rijonx\_scf(...)} &
  A converged pyscfad mean-field object; arguments \texttt{(mol,
  auxbasis=None, rcond, **kernel\_kwargs)}.  \texttt{RIJONX\_ROHF} is
  selected when \texttt{mol.spin} is non-zero, and a non-converged run
  raises. \\
\texttt{rijonx\_scf\_gradient(...)} &
  Full nuclear gradient of the converged closed-shell energy, shape
  \texttt{(natm, 3)} in Eh/bohr; arguments \texttt{(mf, auxbasis,
  rcond)}. \\
\texttt{rijonx\_energy(mol, C, ...)} &
  The energy of Eq.~\eqref{eq:new-rijonx} at fixed coefficients. \\
\texttt{rijonx\_gradient(mol, C, ...)} &
  The fixed-coefficient (skeleton) gradient: RI-J analytic, exchange
  dense, plus core and nuclear terms. \\
\texttt{ri\_j\_*}, \texttt{exact\_k\_*}, \texttt{core\_hamiltonian\_*}, \texttt{nuclear\_repulsion\_gradient} &
  The individual energy and gradient pieces. \\
\texttt{make\_cosx\_grid}, \texttt{cosx\_exchange\_matrix}, \texttt{cosx\_exchange\_energy} &
  A fixed-grid COSX prototype for exchange; not a production route. \\
\bottomrule
\end{tabularx}
\end{center}

The gradient convention is the point that matters.  The functions in
\texttt{driver.ri\_scf} return a fixed-coefficient gradient: the
coefficients are held at their reference values and only the integrals
and the nuclear repulsion move.  That is not the SCF gradient even at a
converged density; the missing orbital-response term reduces to the
overlap-derivative (energy-weighted density) contribution,
$-2\sum_{\mu \in A} Y_{\mu\nu}\,\partial S_{\mu\nu}/\partial R$, with
$Y = 2 C_{\mathrm{occ}}\epsilon C_{\mathrm{occ}}^{T}$.  Without it the
individual terms cancel strongly and the skeleton sum differs from the
SCF gradient by order $1$~Eh/bohr.  \texttt{rijonx\_scf\_gradient}
supplies the completion and is the entry point to use for a converged
SCF.

The auxiliary basis defaults to \texttt{def2-universal-jkfit} and is
resolved by \texttt{make\_auxmol} in \texttt{core.ri}, which falls back
to the orbital basis when a fitting set does not cover every element.
Build a fresh molecule and mean-field object per geometry: the RI factor
cache
and the pyscfad integral caches are keyed on the built molecule and are
not reliably refreshed by an in-place geometry edit.

The exact exchange branch still forms and caches the dense
\texttt{nao}$^{4}$ tensor and, under geometry differentiation, its
dense derivative.  RI-J removes the Coulomb part of that work, not the
exchange part.  The explicit open-shell gradient is not implemented;
\texttt{rijonx\_scf\_gradient} refuses fractional or single
occupations, and the implicit route \texttt{mf.energy\_grad()} is the
only open-shell path and has not been checked against finite
differences.  Range-separated Coulomb is refused because the factors
are built for the full operator.  ECP molecules work when the molecule
is built with an explicit \texttt{ecp} argument.

\begin{lstlisting}[language=Python]
from diffmrc.driver.ri_scf_driver import rijonx_scf, rijonx_scf_gradient

mf = rijonx_scf(mol, auxbasis='def2-universal-jkfit')
e_tot = mf.e_tot

# Closed-shell full gradient (skeleton plus overlap completion):
grad = rijonx_scf_gradient(mf)

# Open shell: only the implicit path is available.
# grad_open = mf.energy_grad().coords
\end{lstlisting}

The energy and gradient are validated on H$_2$O and CO (closed shell),
the OH radical (open-shell energy), and HI/def2-svp with an ECP
(gradient against central differences).  The explicit and implicit
closed-shell gradients agree to $10^{-6}$~Eh/bohr; the measured
RIJONX-minus-dense energy difference on HI is $-2.6\times10^{-5}$~Eh,
which is the J fit, not an implementation error.

\section{Scalar DF-X2C CASSCF}
\label{sec:new-x2c-casscf}

\texttt{core.x2c\_casscf\_df} adds a state-specific and state-averaged
CASSCF layer on the scalar (spin-free) X2C one-electron Hamiltonian
with a density-fitted two-electron part:

\begin{equation}
  E = \sum_{pq}\gamma_{pq} h_{pq}
    + \tfrac{1}{2}\sum_{pqrs}\Gamma_{pqrs}(pq|rs) + E_{\mathrm{nuc}},
  \qquad
  (pq|rs) \approx \sum_{P} w_{pqP} w_{rsP},
  \label{eq:new-x2c}
\end{equation}

with $h$ the spin-free Foldy--Wouthuysen one-electron Hamiltonian,
$w = B S^{-1/2}$ and no picture-change correction on the Coulomb
operator.  The wavefunction is a real spin-adapted determinant
expansion; the orbitals carry no spin-orbit coupling.

\begin{center}
\small
\begin{tabularx}{\textwidth}{>{\raggedright\arraybackslash}p{0.40\textwidth}>{\raggedright\arraybackslash}X}
\toprule
Entry point & Returns and purpose \\
\midrule
\texttt{df\_x2c\_casscf(...)} &
  A \texttt{DFX2CCASSCFResult} with the converged orbitals, the CI
  vectors, the state energies, the state-averaged energy and the
  orbital-gradient norm.  Required arguments \texttt{(mol, ncore, nact,
  nele)}; optional auxiliary basis, \texttt{nroots}, \texttt{weights},
  \texttt{target}, convergence controls, \texttt{theta0} and
  \texttt{rhf}.  \texttt{nroots=1} is the single-state calculation;
  \texttt{weights} default to equal weights. \\
\texttt{df\_x2c\_casscf\_gradient(...)} &
  \texttt{(grad, parts)} for one CI root, shape \texttt{(natm, 3)} in
  Eh/bohr, with the one-electron, two-electron, Pulay and nuclear
  contributions separated.  Arguments \texttt{(mol, result,
  state=None)}; pass the result object, which carries the CAS space
  definition and the factor derivatives. \\
\texttt{df\_x2c\_energies\_at(...)} &
  The state energies at a new geometry with the rotation parameters
  held fixed; arguments \texttt{(mol, result, theta=None)}.  The
  reference orbitals follow the geometry through a displaced SCF and
  are aligned to \texttt{result.c\_ref}. \\
\texttt{df\_x2c\_sa\_state\_gradient(...)} &
  The state-specific gradient of a state-averaged calculation; requires
  \texttt{result.nroots >= 2}.  Arguments \texttt{(mol, result,
  state=None)}, with finite-difference steps \texttt{fd\_step} and
  \texttt{theta\_step}, a pseudo-inverse cutoff \texttt{rcond}, and
  \texttt{return\_info}.  Returns \texttt{(grad, info)} when
  \texttt{return\_info} is true. \\
\texttt{df\_x2c\_rhf(...)} &
  The closed-shell reference SCF on the same Hamiltonian; arguments
  \texttt{(mol, factors, auxmol, auxbasis, conv\_tol, max\_cycle,
  dm0)}. \\
\texttt{sfx2c\_hcore(mol)}, \texttt{spatial\_df\_factors(mol, ...)} &
  The Hamiltonian and the spatial RI factors. \\
\bottomrule
\end{tabularx}
\end{center}

At a converged CASSCF point the orbital and CI responses vanish, so the
single-state gradient is the direct derivative at fixed orbitals and CI
vector: the one-electron term uses the spin-free X2C hcore derivative,
the two-electron term uses the three-center and metric derivatives of
Chapter~\ref{ch:ri-df}, and the Pulay term uses
$-\sum_{pq}F_{pq}\,\partial S_{pq}/\partial R$.

A state-averaged calculation is different: the orbitals minimize
$E_{\mathrm{SA}}$, so the target-state orbital gradient $g_t$ does not
vanish.  The state-specific gradient is

\begin{equation}
  \frac{dE_t}{dR} = \left.\frac{dE_t}{dR}\right|_{\theta}
    + z^{T} M, \qquad H z = -g_t,
    \qquad H = \frac{\partial^2 E_{\mathrm{SA}}}{\partial\theta^{2}},
  \label{eq:new-x2c-sa}
\end{equation}

with $M = \partial g_{\mathrm{SA}}/\partial R$.  Both $H$ and $M$ are
evaluated by finite differences of the analytic state-averaged orbital
gradient: $H$ over $\theta$ (the CI vectors are re-diagonalized at every
step, so the CI response is included) and $M$ over displaced
geometries, rebuilding the integrals and the geometry-dependent
reference orbitals.  The direct term is a central difference of the
target-state energy along the same envelope path; the closed-form
single-state formula does not apply because the target state is not
stationary at the SA orbitals.

The single-state energy and gradient are validated on H$_2$O/6-31G
(6-31G orbital basis, \texttt{def2-universal-jkfit}, CAS(2,2), four
core spatial orbitals) against a density-fitted sfx2c1e CASSCF in
pyscf: the energy agrees within $10^{-5}$~Eh, the gradient within
$10^{-4}$ relative, the envelope finite difference within
$5\times10^{-6}$~Eh/bohr, and the state-averaged state gradients
satisfy the weighted sum rule to $10^{-5}$~Eh/bohr.  The layer is
scalar (spin-free): full-electron basis sets only, because the spin-free
X2C helper raises \texttt{NotImplementedError} for ECP molecules.
Complex two-component CI, Kramers-restricted response and ECP support
are not implemented.  The CI, RDM and Fock contractions are dense in
the spin-orbital space (\texttt{nso}$^{4}$ for the 2-RDM), so the
prototype is sized for small validation systems; production f-block
active spaces need a factorized 2-RDM route.

\begin{lstlisting}[language=Python]
from diffmrc.core import x2c_casscf_df as xc

res = xc.df_x2c_casscf(mol, ncore=4, nact=2, nele=2,
                       auxbasis='def2-universal-jkfit')
grad, parts = xc.df_x2c_casscf_gradient(mol, res)
e_states = xc.df_x2c_energies_at(mol, res)

res_sa = xc.df_x2c_casscf(mol, ncore=4, nact=2, nele=2, nroots=2,
                          weights=[0.5, 0.5],
                          auxbasis='def2-universal-jkfit')
g_state, info = xc.df_x2c_sa_state_gradient(
    mol, res_sa, state=0, return_info=True)
\end{lstlisting}

\section{Validation scope and boundaries}
\label{sec:new-scope}

The table collects the validated systems and the boundaries of the five
layers.  A system absent from the second column is not validated, and
the absence is a statement about coverage, not about expected accuracy.

\begin{center}
\begin{tabularx}{\textwidth}{>{\raggedright\arraybackslash}p{0.19\textwidth}>{\raggedright\arraybackslash}p{0.34\textwidth}X}
\toprule
Module & Validated on & Not covered \\
\midrule
\texttt{nevpt2\_zvector} &
  H$_2$O/6-31G, LiH/6-31G, CAS(2,2), closed shell, frozen core &
  open shell, ECP, f block, large basis; the exact treatment covers the
  \texttt{Si} class only \\
\texttt{caspt2} coupled &
  H$_2$O/6-31G and CO/6-31G, CAS(2,2), against the reference class
  table &
  differentiable coupled path; large active spaces \\
\texttt{somf\_ecp} &
  GeF with a spin-orbit ECP; YbF \texttt{ECPso} norms; operator and
  response identities &
  production absolute SOC; orbital and CI response \\
\texttt{ri\_scf}, \texttt{ri\_scf\_driver} &
  H$_2$O, CO, OH (energy), HI/def2-svp with an ECP (gradient) &
  open-shell analytic gradient; RI-JK; production COSX \\
\texttt{x2c\_casscf\_df} &
  H$_2$O/6-31G, CAS(2,2) and SA(2), against density-fitted sfx2c1e
  CASSCF &
  ECP; complex two-component CI; Kramers response; large active spaces \\
\bottomrule
\end{tabularx}
\end{center}
